\section{Marco conceptual}\label{sec:marco}

\paragraph{ReAct.}
Yao et al.~\cite{react} proponen intercalar trazas de razonamiento (\emph{Thought}) con acciones sobre el entorno (\emph{Action}), cuyo resultado (\emph{Observation}) se agrega al contexto antes del siguiente paso.
Así el modelo puede buscar información, calcular o ejecutar código en vez de depender solo de lo que memorizó.

\paragraph{Recuperación aumentada (RAG).}
Un sistema RAG recupera fragmentos relevantes de una base de conocimiento y los agrega al prompt.
En su forma más simple (\emph{Naive RAG}) hace una sola recuperación y una sola llamada al modelo, sin verificar la respuesta.

\paragraph{Sistemas multi-agente y guardrails.}
Dividir el trabajo entre agentes con roles distintos permite separar \emph{producir} una respuesta de \emph{verificarla}.
Un guardrail es una regla de control que bloquea o devuelve una salida que no cumple ciertas condiciones.
Puede ser determinista, como ejecutar pruebas unitarias, o basarse en otro modelo que juzga la salida.

\paragraph{Evaluación con RAGAS y LLM-as-a-Judge.}
RAGAS~\cite{ragas} define métricas sin referencia para sistemas RAG.
La \emph{Faithfulness} es la fracción de afirmaciones de la respuesta que se pueden inferir del contexto recuperado.
La \emph{Answer Relevance} mide qué tan directamente la respuesta atiende la pregunta.
Ambas se calculan pidiéndole a un LLM que actúe como juez, de modo que heredan sus sesgos y su ruido.
